%% 中文简历模板：请将尖括号中的示例文字替换为你的信息。
\documentclass[zh,compact]{resume}

% 姓名区所占页眉宽度比例；联系信息较多时可适当减小。
\headersplit{0.22}
\fileinfo{\faEdit{} \today}

\name{名字}{姓氏}
\keywords{Go, Python, Backend, Kubernetes, PostgreSQL}
\tagline{求职意向 / 职位名称}

\profile{%
  \mobile{138-0000-0000}
  \email{name@example.com} \\
  \github{your-github}
  \iconlink{\faGlobe}{https://example.com}{example.com} \\
  \university{某某大学}
  \degree{计算机科学与技术 \textbullet\ 本科 \textbullet\ 2022--2026}
}

\begin{document}
\makeheader

%======================================================================
\sectionTitle{个人简介}{\faUser}
%======================================================================
\begin{abstract}
\begin{competences}[4.8em]
  \comptence{求职方向}{后端开发 \textbullet\ 云原生 \textbullet\ 平台工程}
  \comptence{个人概述}{用一至两句话概括你的技术方向、最有代表性的经历和能带来的价值。优先写可验证的成果，例如性能提升、用户规模、交付效率或业务影响。}
\end{competences}
\end{abstract}

%======================================================================
\sectionTitle{教育经历}{\faGraduationCap}
%======================================================================
\begin{educations}
  \education{2022.09}[2026.06]{某某大学}{计算机学院}{计算机科学与技术}{本科}
\end{educations}

%======================================================================
\sectionTitle{实习 / 工作经历}{\faBriefcase}
%======================================================================
\begin{experiences}
  \experience%
    [2025.06]%
    {2025.09}%
    {某某科技有限公司 \aside{后端开发实习生}}%
    [\begin{itemize}
      \item 用动作、技术和结果描述你的贡献，例如：负责订单服务接口开发，将核心接口 P95 延迟降低 30\%。
      \item 用数字说明影响范围，例如日请求量、用户数、成本或故障率变化。
    \end{itemize}]
\end{experiences}

%======================================================================
\sectionTitle{项目经历}{\faCode}
%======================================================================
\begin{projects}
  \project%
    {2025.03}%
    [2025.05]%
    {个人项目}%
    {后端服务}%
    {项目名称}%
    [一句话说明项目解决的问题、目标用户和你的职责。]%
    {\begin{itemize}
      \item \textbf{技术实现}：说明架构或关键技术选择，以及为什么这样做。
      \item \textbf{成果量化}：说明性能、稳定性、效率或用户体验方面的结果。
    \end{itemize}}%
    [Go, Gin, PostgreSQL, Redis, Docker]
\end{projects}

%======================================================================
\sectionTitle{核心技能}{\faWrench}
%======================================================================
\twocolumnsection{%
  \begin{competences}[4.8em]
    \comptence{编程语言}{Go、Python、SQL、Shell}
    \comptence{后端开发}{HTTP / gRPC、REST API、鉴权与测试}
    \comptence{数据存储}{PostgreSQL、MySQL、Redis}
  \end{competences}%
}{%
  \begin{competences}[4.8em]
    \comptence{云原生}{Docker、Kubernetes、CI/CD}
    \comptence{工程能力}{Git、Linux、监控与排障}
    \comptence{其他技能}{按岗位要求补充，不要罗列未掌握的工具}
  \end{competences}%
}

\end{document}
